\section{Desafios e Fronteiras de Escalabilidade}
\label{sec:desafios}

A viabilização do Gêmeo Digital como infraestrutura nativa e onipresente em clínicas e sistemas de saúde baseados em evidência exige a transposição de barreiras regulatórias, computacionais e demográficas extremas.

\subsection{Custo-Efetividade e Complexidade Computacional}
Simulações preditivas que adotam métodos recursivos dinâmicos ou elementos viscoplásticos explodem na casa de dezenas de horas de paralelismo massivo em GPUs avançadas. Ademais, nota-se uma abissal ausência de matrizes econômicas sólidas reportadas \cite{duggal2026exploring}. Para sistemas subsidiados por entidades públicas (como o SUS, Sistema Único de Saúde brasileiro), a equação do impacto em políticas públicas frente à massiva requisição em servidores em nuvem dita a viabilidade em detrimento a tratamentos conservadores diretos.

\subsection{Calibração Fisiológica In Vivo}
As métricas mecânicas que balizam o FEM na literatura são notoriamente variadas. O ligamento periodontal difere mecanicamente dependendo da idade, doenças periodontais e carga histórica de atrito. A injeção generalizada de modelos probabilísticos da coorte sem um mecanismo ágil de parametrização indireta (baseada em exames rotineiros de radiografia e bioimpedância intraoral) contamina os resultados preditivos com erros estruturais ocultos. A modelagem inversa persiste como um funil computacional intensivo a ser mitigado por técnicas de inteligência de máquina eficientes \cite{olawade2026advancing}.

\subsection{Validação Clínica Longitudinal}
Revisões sistemáticas incisivas, como de Duggal et al. (2026), expõem que a ampla maioria dos protótipos de gêmeos digitais jaz em estágios de \textit{prova de conceito} isolados, orbitando majoritariamente ao redor de casos in vitro com amostragens reduzidas. O avanço em direção ao estabelecimento da confiabilidade máxima só ocorrerá quando Ensaios Clínicos Randomizados Controlados (RCTs) multicêntricos conduzirem testes com horizontes superiores a 12 ou 24 meses.

\subsection{Equidade Algorítmica e Desertos Digitais}
A inserção indiscriminada de complexos digitais no \textit{core} da odontologia pode engendrar exclusão sociomédica violenta, marginalizando populações situadas em \textit{desertos digitais} \cite{katsoulakis2024digital}. Uma integração tele-odontológica estruturada deve transacionar predições de baixa latência em interfaces leves (através de compactação e métodos de ordem reduzida - ROM) a agentes e profissionais operando na extremidade, mitigando vieses algorítmicos em bancos de dados etnicamente concentrados.

\subsection{Normatização e Interoperabilidade FHIR}
Padrões de governança unificados como a ISO/IEC 3173 e ecossistemas open-source baseados no padrão HL7 FHIR devem balizar as linguagens comuns entre laboratórios parceiros, hospitais oncológicos maxilofaciais e clínicas odontológicas. O encapsulamento ético dos dados deve mitigar as vulnerabilidades persistentes à privacidade inerentes à cópia computacional de perfis biométricos altamente identificáveis e invioláveis.
